
\begin{theorem}[Theorem 18.5]\label{mcrt-auto-theorem-18-5}\dcref{18.5}\source{matsumura-commutative-ring-theory:page-0161}\uses{}
\begin{mathematicalrecord}
We consider modules over a Noetherian ring A.
  (i) A direct sum of any number of injective modules is injective.
  (ii) Every injective module is a direct sum of indecomposable injective
modules.
  (iii) The direct sum decomposition in (ii) is unique, in the sense that if
         M = @ Mi (with indecomposable          Mi)
then for any PESpec A, the sum M(P) of all the Mi isomorphic to E(A/P)
depends only on M and P, and not on the decomposition M = @ Mi.
Moreover, the number of Mi isomorphic to E(A/P) is equal to
          dim,cp,Hom,,(rc(P), MJ, (where K(P) = AdPA,),
so that this also is independent of the decomposition.
\end{mathematicalrecord}
\end{theorem}
\begin{proof}
\begin{mathematicalrecord}
(i) Let M1 for 1~11 be injective modules. It is enough to prove that
for an ideal I of A, any linear map q : I -+ @MI can be extended to a linear
e I is finitely generated, q(Z) is contained
                                             roftheM,.Ifcp(l)cM,@...@M,and
                                            t of p(a) in Mi, then Cpi:I -+ Mi extends
                                            -@J%             by t41)=11/,(1)+.~.+$,(1)

b. (ii) Say that a family .F = {E} of indecomposable injective submodules
& of M is fvee if the sum in M of the E, is direct, that is if, for any finite
$ number EA,,. . , E,,, of them,
$;
./I           E,, n(E,,    + ... + E,?) = 0.
&tit        %JIbe the set of all free families 8, ordered by inclusion. Then by
2: Zorn?s lemma !JJI has a maximal element, say FO. Write N = CEEFOE; then
g by (i), N is injective, hence a direct summand of M, and M = N 0 N?. If
g, NVf 0 then since it is a direct summand of M it must be injective, and
wa.for PEAss(N?), the proof of Theorem 4, (ii), shows that N? contains a
     direct summand E? isomorphic to E(A/P). Thus F0 u {E?} is a free family,
     contradicting the maximality of FO. Hence N? = 0 and M = N.
        (iii) If we can show that M(P) has the property that every submodule
     & of M isomorphic to E(A/P) is contained in M(P), then M(P) is the
     submodule of M generated by all such E, and therefore is determined by
   ?M and P only. To prove this, take any <EE; we can write t = 4, + ... +
     4, with ~iEM(Pi), where PI,. . . , P, are distinct prime ideals and P = P,.
     Setting~,-5=vl,and5i=rifor2~idrwehaver],+...+vl,=O,with
     )1IcM(P,) + E and qieM(Pi) for i >, 2. We need only prove that in this case
     each Y/i= 0. Suppose that P, is minimal among P,, . . , P,; then for any m
     we have (PI . . .Prml)??@Pl,        so that taking ag(P,...P,-,)?-P,      and m
     large enough, we get ar], = ... = ar],- 1 = 0. Then also q, = 0, but
     multiplication     by a is an automorphism         of M(P,), so that ql = 0. By
L ?mdnction on r we get vi = 0 for all i.
f       We now prove that if M(P) = M, @... @ M, with Mi N E(A/P) then
$                s = dim,(,,) Hom,,(rc(P), Mp).
I;
$ (we are writing this as ifs were finite, but, as one can see from the proof
k; below, the same works for any cardinal number.) By Theorem 4, (vi), both
     aides of M(P) = M, @...@ M, are &modules,                       and Mi = E(rc(P)).
     Moreover, by Theorem 4, (v), E(LI/Q)~ = 0 if Q $ P, so that

              MP = @ M(Q)p = @ M(Q).
                     QcP             QcP

      Hence we can replace A by A p, and assume that A is a local ring with P
      its maximal ideal; set k = K(P). If Q # P then any XEP - Q gives an
      automorphism of M(Q), but x.k = 0, so that Hom,(k, M(Q)) = 0. Hence
      Hom,(k, M) = Hom,(k, M(P)), so that there is no loss of generality in
assuming that M = M(P). For any A-module N, we can identify,
Hom,(k,N)     with the submodule {<ENJP~ =O}, but since E(k) is an
essential extension of k we must have dim, Hom,(k, E(k)) = 1, so that if
M = M, @...@ M, with Mi N E(k) then s = dim, Hom,(k, M). H
\end{mathematicalrecord}
\end{proof}


\begin{theorem}[Theorem 18.6]\label{mcrt-auto-theorem-18-6}\dcref{18.6}\source{matsumura-commutative-ring-theory:page-0163}\uses{mcrt-auto-theorem-3-1}
\begin{mathematicalrecord}
Let (A, m, k) be a Noetherian local ring, and E = E,(k) the
injective hull of k. For each A-module M set M? = Hom,(M, E).
    (i) If M is an A-module and 0 # XEM, then there exists cp~M? such that
q(x) # 0. In other words the canonical map 8: M --+ M? defined by Q)(,)
 = q(x) for XEM and cp~M? is injective.
    (ii) If M is an A-module of finite length, then l(M) = 1(M?) and the
canonical map M -M?            is an isomorphism.
    (iii) Let A^ be the completion of A; then E is also an A-module, and
is an injective hull of k as A-module.
    (iv) Hom,(E, E) = Hom,(E, E) = A^. In other words, each endomer-
phism of the A-module E is multiplication         by a unique element of 2.
    (v) E is Artinian as an A-module and also as an A-module. Assume
now that A is complete, and write N(resp. ~2) for the category of
Noetherian (respectively Artinian) A-modules. Then if ME&? we have
M?E~       and MN M?; if ME& we have M?EJ+?? and MN M?.
\end{mathematicalrecord}
\end{theorem}
\begin{proof}
\begin{mathematicalrecord}
(i) Let f:Ax -E           be the composite of the canonical maps
Ax N A/ann(x), A/ann(x) - A/m = k and k - E. Then f(x) # 0. Since Eis
injective we can extend f to cp:M -E.
   (ii) If l(M) = n < cc then M has a submodule M, of length n - 1, and
O-+M,-M-k-,Oisexact,sothatO-+k?--+M?--rM~+Oisexact.
However
            k? = Horn (k, E) = Horn (k, k) ru k,
so that by induction on II we get l(M) = n = [(M?). The canonical map
M -M?        is injective by (i), and I(M) = QM?) = l(M?), hence it must be an
isomorphism.
   (iii) Each element of E is annihilated by some power of m, so that the
canonical map E -E         @AA is surjective. However, since A^ is faithfully
flat over A it is also injective, so that E N E aAA, and we can view E as
an A-module. Let F be the injective hull of E as an A-module. Then F
is also the A-injective hull of k, so that every element of F is annihilated
by some power of rn2. As an A-module F splits into a direct sum Of E
and an A-module C. If xgC, and if mrAx = 0, then for each a*Ea we
can find aeA such that a* -a mod m?A and hence a*x =ax EC.
Therefore C is an A-module. But F is indecomposable as an A-module?
Hence C = 0 and E = F.
    (iv) For v > 0 set E, = (xEE)m?x = O}. Then we have (A/m?yS
Horn, (A/m?, E) rr E,, and Horn, (E,, E,) = Hom,(E,, E) = E:. = (Ah?)? =
A/m. Now El = E, = . . . and E = UYEY by Theorem 4, (v), hence E =
limE,. Therefore Hom,(E, E) = Horn,@     E,, E) = @ HomA(E,, E) =
GA/m?     = A^.
7?) Ifan A-module M is Artinian and XEM, then Ax in A/ann(x) is also
 Artinian and consequently my c arm(x) for some v. Therefore M can be
viewed as an A-module, and its A-submodules are precisely its A-
submodules. It is also clear that if an a-module M is Artinian then we have
the same conclusion. Therefore to prove (v) we may assume that A is
complete.
   JfM is a submodule of E set ML = (UEA jaM = O}. If Z is an ideal of A set
Zl= {x~EjZx =O}. Then clearly ML? 3 M. If XEE - M there exist
&E/M?)        satisfying cp(xmod M) # 0 by (i), and if we identify E? =
Hom,(E, E) with A then (E/M)? is identified with Ml. Thus cp(x mod M) =
ax for some aeM?,         and .x$M?.      Therefore Ml? = M. Similarly, if
acgA -I     then there exists cpe(A/I)? such that cp(umodZ) #O, and
(A/Z)? is identified with the submodule I1 of E = A?. Thus, setting x =
cp(1modZ) we have xel? and ux = cp(umod I) # 0. This proves a$Z?, so
that I = Z1l. Thus M w ML is an order-reversing bijection from the set of
submodules of E onto the set of ideals of A. Since A is Noetherian, it follows
that E is Artinian. By Theorem 3.1 finite direct sums E? of E are also
Artinian for all n > 0.
   If MEN then there is a surjection A? + M for some n, and so there
is an injection M? -(A?)?      = E?. Hence M? is Artinian. On the other
 hand, if ME& there is an injection M -En           for some n. This can be
seen as follows: consider all linear maps M + E?, where n is not fixed,
and take one cp:M -En         whose kernel is minimal among the kernels of
 those maps. Then, using (i) we can easily see that Ker(cp) = 0. Now, from
0 + M - E? we have (E?)? = A? - M? + 0 exact, hence M?eJlr. Now the
 assertion M z M? for M EJY or d can easily be checked using (iv) if M = A
 or E, and the general case follows from this and from (i). n
Zhnma 6. Let A be a Noetherian ring, S c A a multiplicative      set, M an A-
module and N c M a submodule. Assume that M is an essential extension
of N; then MS is an essential extension of N,.
proof. For (EM we write ts = l/l EMU; then any element of MS can be
written u*cs (with u a unit of A, and REM), so that it is enough to show that
for any non-zero & we have N,n A,.<, # 0. Suppose that ann (tot) is a
maximal element of the set of ideals {ann(    tES}; then ifwe set V= to 4, we
have ts = to 1qs, and hence v ~0. Now let b = {~EAI~?EN};                  by
assumption,
         bv=AqnN#O.
Suppose that b = (b ?,...,b,);   if b,q,=...   = b,qs =0 then there is a tcS
150      Regular sequences

such that tb,r = 0 for all i. Then tbv] = 0, but by choice of q we have
ann (q) = ann (tq), so that bq = 0, which is a contradiction. Thus biqs # 0
for some i, and
           biYls~As~V,nNs = As~s~Ns,
 as required.     n
    By Lemmas 5 and 6, if M is an injective hull of N then the As-module
M, is an injective hull of N,. Hence if O+ M - I0 -+I?          - ... is
a minimal injective resolution of an A-module M, then 0 + M, --+ I,0 -+
G -...      is a minimal injective resolution of the As-module M,. The
I? are determined uniquely up to isomorphism by M. We can therefore
define pi(P, M) to be the number of summands isomorphic to E(A/P)
appearing in a decomposition of I? as a direct sum of indecomposable
modules. We can write symbolically
         I?=          @       pi(P, M)E(A/P).
                  Pdpec   A

From what we have just proved, for a multiplicative          set S c A,
      Ili(P, M) = ~i(PA,, MJ   if P n S = a.
\end{mathematicalrecord}
\end{proof}


\begin{theorem}[Theorem 18.7]\label{mcrt-auto-theorem-18-7}\dcref{18.7}\source{matsumura-commutative-ring-theory:page-0165}\uses{}
\begin{mathematicalrecord}
Let A be a Noetherian                 ring, M an A-module, and
PESpec A. Then
           Pi(P, W = dimKcp, Ext&(lc(P), MP) = dim,&Ext>(A/P,          M))p.
In particular, if M is a finite A-module then ~i(P, M) < co.
\end{mathematicalrecord}
\end{theorem}
\begin{proof}
\begin{mathematicalrecord}
Replacing A and M by A, and M, we can assume that (A, P, k) is
a local ring. Let O+M -+ZO-f+ll--%...                     be a minimal    injective
resolution of M, so that Exti(k, M) is obtained as the homology of the
complex
          . ..-+Hom.(k,I?-?)-Hom,(k,Z?)-Hom,(k,I?+?)
               +...
We can identify Hom,(k,I?)      with the submodule T?= {x~Z?/Px =0} c
I?. By construction of the minimal injective resolution, I? is an essential
extension of d(l?- ?), so that for XGT? the submodule AX 1: k intersects
d(l?-?), and xEd(Z?-?). Therefore, T? c d(I?-?), and dT?-? = dT?= 0, so
that Ext$(k, A) = T?. Also,
         dim, T? = dim, Hom,(k, I?),
and by Theorem 4, (iii), this is equal to pi(P, M). w
\end{mathematicalrecord}
\end{proof}


\begin{theorem}[Theorem 18.8]\label{mcrt-auto-theorem-18-8}\dcref{18.8}\source{matsumura-commutative-ring-theory:page-0165}\uses{}
\begin{mathematicalrecord}
A necessary and sufficient condition for a ring A to be
Gorenstein is that a minimal     injective resolution O+ A -+ 1? -+
I? --+ .. . of A satisfies
         I? = @           E(A/P),
                 htP=i
or, in other words, /li(P, A) = Si,htP (the Kronecker 6) for every PESpec A.
\end{mathematicalrecord}
\end{theorem}
\begin{proof}
\begin{mathematicalrecord}
By Theorem 7 and condition (2) of Theorem 1 we have
                A, is Gorenstein e pi(P, A) = 8i,htP.
\end{mathematicalrecord}
\end{proof}


\begin{theorem}[Theorem 18.9]\label{mcrt-auto-theorem-18-9}\dcref{18.9}\source{matsumura-commutative-ring-theory:page-0166}\uses{}
\begin{mathematicalrecord}
Let (A,m)       be a Noetherian     local ring, and M a finite
      A-module. Then
              inj dim M < r3 +inj dim M = depth A.
\end{mathematicalrecord}
\end{theorem}
\begin{proof}
\begin{mathematicalrecord}
Suppose that inj dim M = r < cc. If P is a prime ideal distinct from
      m, choose xEm - P. Then

ii            0 + A/P 2      A/P,
      together with the right-exactness      of Ext>( - , M) gives an exact sequence
              Ext?,(A/P,M)AExt:,(A/P,M)+O,
)..
(i.   so that by NAK Ext>(A/P, M) = 0. Putting this together with Lemma 1,
:.-
 :    we get Ext>(k,M)#O.             Set t=depthA,   and let xr,,..,x,Em        be a
$     maximal A-sequence; then setting A/(x,, . ,x,) = N we have mEAss (N).
$     Hence there exists an exact sequence O-+ k - N, and we must have
!-    Ext>(N, M) # 0. The Koszul complex K(x r, . . . ,x,) is a projective resolution
?1
2:
~:-   of N = A/(x,, . . . , x,), so computing Ext by means of it we see that
              Ext;(N, M) N M/(x,, . . ,x,)M,
;.,,.
$, and by NAK this is non-zero. Thus proj dim N = t, and from Ext\(N, M) #
?? 0 we get t d r, whereas from Ext>(N, M) # 0 we get t 2 r. Hence t = r. H
8 Remark (the Bass conjecture and the intersection theorem). Let (A, m, k) be
i,
; a Noetherian local ring of dimension d. H. Bass [l] conjectured the
1. following:
[     (B)     if there exists a finite A-module         M (#Q     of finite   injective
8:?           dimension, then A is a CM ring.
F According to Theorem 9 this is equivalent to asking that inj dim M = d.
i The converse of the Bass conjecture is true. Indeed, if A is CM, taking a
d maximal A-sequence x1,. . , x,, and setting B = A/(x,, . . . , xd) and E = E(k)
$?
!+ we have l,(B) < co. By Theorem 6, M = Horn,@, E) is also of finite length,
in hence is finitely generated. We prove inj dim,M < d; the Koszul complex
? O--bA-+Ad -+~~~-+Ad--tA+BL+O with respect to x1,. . . , x, provides an A-
! free resolution of B. Now applying the exact functor Hom,(-,E) to this
    gives    the exact sequence O--+M-+E+Ed-+..~+Ed--+E+O.            This proves
    injdim,M d d.
i.
L        (B) is a special case of the following theorem.
1 (C)         If I?: 0-10 -+?..-+Zd+O is a complex of injective modules such
c             that #(I?) is finitely generated for all i and I? is not exact, then
              Id # 0.
c4
Using the theory of dualizing complexes (see [Rob])            one can prove that (C)
is equivalent to the following
Intersection Theorem. If F.: O+F,-+~~~-+F,-+O         is a complex of finitely
generated free modules such that H,(F) has finite length for all i and F. is
not exact, then F, # 0.
   (B) was proved by Peskine and Szpiro [l] in some important cases,and
by Hochster [H] in the equal characteristic case (i.e. when A contains a
field) as a corollary of his existence theorem for the ?big CM module?, i.e. a
(not necessarily finite) A-module with depth = dim A, see [H] p.10 and
p.70. The intersection theorem was conjectured by Peskine-Szpiro [3] and
by P. Roberts independently. They pointed out that it was also a
consequence of Hochster?s theorem. Finally, P. Roberts [3] settled the
remaining unequal characteristic case of the intersection theorem by using
the advanced technique of algebraic geometry developed by W. Fulton
([Full). Therefore (B), which was known as Bass?sconjecture for 24 years,
is now a theorem. Some other conjectures listed in [H] are still open.
\end{mathematicalrecord}
\end{proof}
